% Conceptual mechanism diagram: IVF operator-host compatibility hypothesis.
%
% Usage in main tex file:
%   \begin{figure}[t]
%     \centering
%     % Conceptual mechanism diagram: IVF operator-host compatibility hypothesis.
%
% Usage in main tex file:
%   \begin{figure}[t]
%     \centering
%     % Conceptual mechanism diagram: IVF operator-host compatibility hypothesis.
%
% Usage in main tex file:
%   \begin{figure}[t]
%     \centering
%     % Conceptual mechanism diagram: IVF operator-host compatibility hypothesis.
%
% Usage in main tex file:
%   \begin{figure}[t]
%     \centering
%     \input{figures/hosts_v2_fig9_mechanism}
%     \caption{Causal hypothesis for IVF host compatibility. Hosts with directional
%              selection pressure (SPEA2, NSGA-III) create feedback that guides IVF
%              intensification toward under-covered regions, yielding strong synergy.
%              NSGA-II's locally isotropic crowding distance provides no such signal,
%              leaving IVF without directional guidance and yielding weak synergy.}
%     \label{fig:mechanism}
%   \end{figure}
%
% Required in preamble:
%   \usepackage{tikz}
%   \usetikzlibrary{positioning,calc,arrows.meta}

\begin{tikzpicture}[
  % --- global style ---
  node distance = 1.4cm and 1.8cm,
  font          = \small,
  >=stealth,
  % node styles
  host/.style   = {
    rounded corners=4pt, draw, thick,
    minimum width=2.6cm, minimum height=0.7cm,
    align=center
  },
  stage/.style  = {
    rounded corners=4pt, draw, thick,
    minimum width=3.0cm, minimum height=0.7cm,
    align=center
  },
  outcome/.style = {
    rounded corners=4pt, draw, thick,
    minimum width=2.4cm, minimum height=0.7cm,
    align=center
  },
  ivfnode/.style = {
    rounded corners=4pt, draw=orange!70!black, thick,
    fill=yellow!25, minimum width=2.6cm, minimum height=0.7cm,
    align=center
  },
  % arrow styles
  fwd/.style    = {->, thick},
  dashed_arr/.style = {->, dashed, thick, orange!70!black},
]

%% ---------------------------------------------------------------
%% Row 1: SPEA2 / NSGA-III path (blue tones)
%% ---------------------------------------------------------------

\node[host, fill=blue!15, draw=blue!60]
  (hostA) {SPEA2 / NSGA-III};

\node[stage, fill=blue!10, draw=blue!50, right=of hostA]
  (pressA) {Directional\\selection pressure};

\node[stage, fill=blue!10, draw=blue!50, right=of pressA]
  (signalA) {Feedback on\\under-covered regions};

\node[outcome, fill=green!20, draw=green!60!black, right=of signalA]
  (outA) {\textbf{Strong synergy}};

%% ---------------------------------------------------------------
%% Row 2: NSGA-II path (red tones)
%% ---------------------------------------------------------------

\node[host, fill=red!15, draw=red!60, below=2.4cm of hostA]
  (hostB) {NSGA-II};

\node[stage, fill=red!10, draw=red!50, right=of hostB]
  (pressB) {Local isotropic\\diversity pressure};

\node[stage, fill=red!10, draw=red!50, right=of pressB]
  (signalB) {No directional\\signal to IVF};

\node[outcome, fill=orange!25, draw=orange!70!black, right=of signalB]
  (outB) {\textbf{Weak synergy}};

%% ---------------------------------------------------------------
%% Shared IVF operator node (centered between the two rows)
%% ---------------------------------------------------------------

% Place it vertically between the two paths, aligned with pressA/pressB
\node[ivfnode, below=1.0cm of pressA, yshift=0.3cm]
  (ivf) {IVF Operator};
% Manually reposition to sit exactly halfway
\coordinate (midAB) at ($(pressA)!0.5!(pressB)$);
\node[ivfnode] (ivf) at (midAB) {IVF Operator};

%% ---------------------------------------------------------------
%% Arrows — top path
%% ---------------------------------------------------------------

\draw[fwd, blue!70]  (hostA)   -- (pressA);
\draw[fwd, blue!70]  (pressA)  -- (signalA);
\draw[fwd, green!60!black] (signalA) -- (outA)
  node[midway, above, font=\scriptsize, text=green!50!black]
  {IVF aligns};

%% ---------------------------------------------------------------
%% Arrows — bottom path
%% ---------------------------------------------------------------

\draw[fwd, red!70]   (hostB)   -- (pressB);
\draw[fwd, red!70]   (pressB)  -- (signalB);
\draw[fwd, orange!80!black] (signalB) -- (outB)
  node[midway, above, font=\scriptsize, text=orange!60!black]
  {IVF drifts};

%% ---------------------------------------------------------------
%% Dashed arrows from IVF node to both paths
%% ---------------------------------------------------------------

\draw[dashed_arr] (ivf.north) -- (pressA.south)
  node[midway, left, font=\scriptsize, text=orange!70!black]
  {intensifies};

\draw[dashed_arr] (ivf.south) -- (pressB.north)
  node[midway, left, font=\scriptsize, text=orange!70!black]
  {intensifies};

%% ---------------------------------------------------------------
%% Labels at far left (hypothesis annotation)
%% ---------------------------------------------------------------

\node[font=\scriptsize\itshape, text=blue!70, left=0.5cm of hostA,
      align=right, anchor=east]
  {\textbf{H:} directional\\feedback};

\node[font=\scriptsize\itshape, text=red!70, left=0.5cm of hostB,
      align=right, anchor=east]
  {\textbf{H:} isotropic\\feedback};

\end{tikzpicture}

%     \caption{Causal hypothesis for IVF host compatibility. Hosts with directional
%              selection pressure (SPEA2, NSGA-III) create feedback that guides IVF
%              intensification toward under-covered regions, yielding strong synergy.
%              NSGA-II's locally isotropic crowding distance provides no such signal,
%              leaving IVF without directional guidance and yielding weak synergy.}
%     \label{fig:mechanism}
%   \end{figure}
%
% Required in preamble:
%   \usepackage{tikz}
%   \usetikzlibrary{positioning,calc,arrows.meta}

\begin{tikzpicture}[
  % --- global style ---
  node distance = 1.4cm and 1.8cm,
  font          = \small,
  >=stealth,
  % node styles
  host/.style   = {
    rounded corners=4pt, draw, thick,
    minimum width=2.6cm, minimum height=0.7cm,
    align=center
  },
  stage/.style  = {
    rounded corners=4pt, draw, thick,
    minimum width=3.0cm, minimum height=0.7cm,
    align=center
  },
  outcome/.style = {
    rounded corners=4pt, draw, thick,
    minimum width=2.4cm, minimum height=0.7cm,
    align=center
  },
  ivfnode/.style = {
    rounded corners=4pt, draw=orange!70!black, thick,
    fill=yellow!25, minimum width=2.6cm, minimum height=0.7cm,
    align=center
  },
  % arrow styles
  fwd/.style    = {->, thick},
  dashed_arr/.style = {->, dashed, thick, orange!70!black},
]

%% ---------------------------------------------------------------
%% Row 1: SPEA2 / NSGA-III path (blue tones)
%% ---------------------------------------------------------------

\node[host, fill=blue!15, draw=blue!60]
  (hostA) {SPEA2 / NSGA-III};

\node[stage, fill=blue!10, draw=blue!50, right=of hostA]
  (pressA) {Directional\\selection pressure};

\node[stage, fill=blue!10, draw=blue!50, right=of pressA]
  (signalA) {Feedback on\\under-covered regions};

\node[outcome, fill=green!20, draw=green!60!black, right=of signalA]
  (outA) {\textbf{Strong synergy}};

%% ---------------------------------------------------------------
%% Row 2: NSGA-II path (red tones)
%% ---------------------------------------------------------------

\node[host, fill=red!15, draw=red!60, below=2.4cm of hostA]
  (hostB) {NSGA-II};

\node[stage, fill=red!10, draw=red!50, right=of hostB]
  (pressB) {Local isotropic\\diversity pressure};

\node[stage, fill=red!10, draw=red!50, right=of pressB]
  (signalB) {No directional\\signal to IVF};

\node[outcome, fill=orange!25, draw=orange!70!black, right=of signalB]
  (outB) {\textbf{Weak synergy}};

%% ---------------------------------------------------------------
%% Shared IVF operator node (centered between the two rows)
%% ---------------------------------------------------------------

% Place it vertically between the two paths, aligned with pressA/pressB
\node[ivfnode, below=1.0cm of pressA, yshift=0.3cm]
  (ivf) {IVF Operator};
% Manually reposition to sit exactly halfway
\coordinate (midAB) at ($(pressA)!0.5!(pressB)$);
\node[ivfnode] (ivf) at (midAB) {IVF Operator};

%% ---------------------------------------------------------------
%% Arrows — top path
%% ---------------------------------------------------------------

\draw[fwd, blue!70]  (hostA)   -- (pressA);
\draw[fwd, blue!70]  (pressA)  -- (signalA);
\draw[fwd, green!60!black] (signalA) -- (outA)
  node[midway, above, font=\scriptsize, text=green!50!black]
  {IVF aligns};

%% ---------------------------------------------------------------
%% Arrows — bottom path
%% ---------------------------------------------------------------

\draw[fwd, red!70]   (hostB)   -- (pressB);
\draw[fwd, red!70]   (pressB)  -- (signalB);
\draw[fwd, orange!80!black] (signalB) -- (outB)
  node[midway, above, font=\scriptsize, text=orange!60!black]
  {IVF drifts};

%% ---------------------------------------------------------------
%% Dashed arrows from IVF node to both paths
%% ---------------------------------------------------------------

\draw[dashed_arr] (ivf.north) -- (pressA.south)
  node[midway, left, font=\scriptsize, text=orange!70!black]
  {intensifies};

\draw[dashed_arr] (ivf.south) -- (pressB.north)
  node[midway, left, font=\scriptsize, text=orange!70!black]
  {intensifies};

%% ---------------------------------------------------------------
%% Labels at far left (hypothesis annotation)
%% ---------------------------------------------------------------

\node[font=\scriptsize\itshape, text=blue!70, left=0.5cm of hostA,
      align=right, anchor=east]
  {\textbf{H:} directional\\feedback};

\node[font=\scriptsize\itshape, text=red!70, left=0.5cm of hostB,
      align=right, anchor=east]
  {\textbf{H:} isotropic\\feedback};

\end{tikzpicture}

%     \caption{Causal hypothesis for IVF host compatibility. Hosts with directional
%              selection pressure (SPEA2, NSGA-III) create feedback that guides IVF
%              intensification toward under-covered regions, yielding strong synergy.
%              NSGA-II's locally isotropic crowding distance provides no such signal,
%              leaving IVF without directional guidance and yielding weak synergy.}
%     \label{fig:mechanism}
%   \end{figure}
%
% Required in preamble:
%   \usepackage{tikz}
%   \usetikzlibrary{positioning,calc,arrows.meta}

\begin{tikzpicture}[
  % --- global style ---
  node distance = 1.4cm and 1.8cm,
  font          = \small,
  >=stealth,
  % node styles
  host/.style   = {
    rounded corners=4pt, draw, thick,
    minimum width=2.6cm, minimum height=0.7cm,
    align=center
  },
  stage/.style  = {
    rounded corners=4pt, draw, thick,
    minimum width=3.0cm, minimum height=0.7cm,
    align=center
  },
  outcome/.style = {
    rounded corners=4pt, draw, thick,
    minimum width=2.4cm, minimum height=0.7cm,
    align=center
  },
  ivfnode/.style = {
    rounded corners=4pt, draw=orange!70!black, thick,
    fill=yellow!25, minimum width=2.6cm, minimum height=0.7cm,
    align=center
  },
  % arrow styles
  fwd/.style    = {->, thick},
  dashed_arr/.style = {->, dashed, thick, orange!70!black},
]

%% ---------------------------------------------------------------
%% Row 1: SPEA2 / NSGA-III path (blue tones)
%% ---------------------------------------------------------------

\node[host, fill=blue!15, draw=blue!60]
  (hostA) {SPEA2 / NSGA-III};

\node[stage, fill=blue!10, draw=blue!50, right=of hostA]
  (pressA) {Directional\\selection pressure};

\node[stage, fill=blue!10, draw=blue!50, right=of pressA]
  (signalA) {Feedback on\\under-covered regions};

\node[outcome, fill=green!20, draw=green!60!black, right=of signalA]
  (outA) {\textbf{Strong synergy}};

%% ---------------------------------------------------------------
%% Row 2: NSGA-II path (red tones)
%% ---------------------------------------------------------------

\node[host, fill=red!15, draw=red!60, below=2.4cm of hostA]
  (hostB) {NSGA-II};

\node[stage, fill=red!10, draw=red!50, right=of hostB]
  (pressB) {Local isotropic\\diversity pressure};

\node[stage, fill=red!10, draw=red!50, right=of pressB]
  (signalB) {No directional\\signal to IVF};

\node[outcome, fill=orange!25, draw=orange!70!black, right=of signalB]
  (outB) {\textbf{Weak synergy}};

%% ---------------------------------------------------------------
%% Shared IVF operator node (centered between the two rows)
%% ---------------------------------------------------------------

% Place it vertically between the two paths, aligned with pressA/pressB
\node[ivfnode, below=1.0cm of pressA, yshift=0.3cm]
  (ivf) {IVF Operator};
% Manually reposition to sit exactly halfway
\coordinate (midAB) at ($(pressA)!0.5!(pressB)$);
\node[ivfnode] (ivf) at (midAB) {IVF Operator};

%% ---------------------------------------------------------------
%% Arrows — top path
%% ---------------------------------------------------------------

\draw[fwd, blue!70]  (hostA)   -- (pressA);
\draw[fwd, blue!70]  (pressA)  -- (signalA);
\draw[fwd, green!60!black] (signalA) -- (outA)
  node[midway, above, font=\scriptsize, text=green!50!black]
  {IVF aligns};

%% ---------------------------------------------------------------
%% Arrows — bottom path
%% ---------------------------------------------------------------

\draw[fwd, red!70]   (hostB)   -- (pressB);
\draw[fwd, red!70]   (pressB)  -- (signalB);
\draw[fwd, orange!80!black] (signalB) -- (outB)
  node[midway, above, font=\scriptsize, text=orange!60!black]
  {IVF drifts};

%% ---------------------------------------------------------------
%% Dashed arrows from IVF node to both paths
%% ---------------------------------------------------------------

\draw[dashed_arr] (ivf.north) -- (pressA.south)
  node[midway, left, font=\scriptsize, text=orange!70!black]
  {intensifies};

\draw[dashed_arr] (ivf.south) -- (pressB.north)
  node[midway, left, font=\scriptsize, text=orange!70!black]
  {intensifies};

%% ---------------------------------------------------------------
%% Labels at far left (hypothesis annotation)
%% ---------------------------------------------------------------

\node[font=\scriptsize\itshape, text=blue!70, left=0.5cm of hostA,
      align=right, anchor=east]
  {\textbf{H:} directional\\feedback};

\node[font=\scriptsize\itshape, text=red!70, left=0.5cm of hostB,
      align=right, anchor=east]
  {\textbf{H:} isotropic\\feedback};

\end{tikzpicture}

%     \caption{Causal hypothesis for IVF host compatibility. Hosts with directional
%              selection pressure (SPEA2, NSGA-III) create feedback that guides IVF
%              intensification toward under-covered regions, yielding strong synergy.
%              NSGA-II's locally isotropic crowding distance provides no such signal,
%              leaving IVF without directional guidance and yielding weak synergy.}
%     \label{fig:mechanism}
%   \end{figure}
%
% Required in preamble:
%   \usepackage{tikz}
%   \usetikzlibrary{positioning,calc,arrows.meta}

\begin{tikzpicture}[
  % --- global style ---
  node distance = 1.4cm and 1.8cm,
  font          = \small,
  >=stealth,
  % node styles
  host/.style   = {
    rounded corners=4pt, draw, thick,
    minimum width=2.6cm, minimum height=0.7cm,
    align=center
  },
  stage/.style  = {
    rounded corners=4pt, draw, thick,
    minimum width=3.0cm, minimum height=0.7cm,
    align=center
  },
  outcome/.style = {
    rounded corners=4pt, draw, thick,
    minimum width=2.4cm, minimum height=0.7cm,
    align=center
  },
  ivfnode/.style = {
    rounded corners=4pt, draw=orange!70!black, thick,
    fill=yellow!25, minimum width=2.6cm, minimum height=0.7cm,
    align=center
  },
  % arrow styles
  fwd/.style    = {->, thick},
  dashed_arr/.style = {->, dashed, thick, orange!70!black},
]

%% ---------------------------------------------------------------
%% Row 1: SPEA2 / NSGA-III path (blue tones)
%% ---------------------------------------------------------------

\node[host, fill=blue!15, draw=blue!60]
  (hostA) {SPEA2 / NSGA-III};

\node[stage, fill=blue!10, draw=blue!50, right=of hostA]
  (pressA) {Directional\\selection pressure};

\node[stage, fill=blue!10, draw=blue!50, right=of pressA]
  (signalA) {Feedback on\\under-covered regions};

\node[outcome, fill=green!20, draw=green!60!black, right=of signalA]
  (outA) {\textbf{Strong synergy}};

%% ---------------------------------------------------------------
%% Row 2: NSGA-II path (red tones)
%% ---------------------------------------------------------------

\node[host, fill=red!15, draw=red!60, below=2.4cm of hostA]
  (hostB) {NSGA-II};

\node[stage, fill=red!10, draw=red!50, right=of hostB]
  (pressB) {Local isotropic\\diversity pressure};

\node[stage, fill=red!10, draw=red!50, right=of pressB]
  (signalB) {No directional\\signal to IVF};

\node[outcome, fill=orange!25, draw=orange!70!black, right=of signalB]
  (outB) {\textbf{Weak synergy}};

%% ---------------------------------------------------------------
%% Shared IVF operator node (centered between the two rows)
%% ---------------------------------------------------------------

% Place it vertically between the two paths, aligned with pressA/pressB
\node[ivfnode, below=1.0cm of pressA, yshift=0.3cm]
  (ivf) {IVF Operator};
% Manually reposition to sit exactly halfway
\coordinate (midAB) at ($(pressA)!0.5!(pressB)$);
\node[ivfnode] (ivf) at (midAB) {IVF Operator};

%% ---------------------------------------------------------------
%% Arrows — top path
%% ---------------------------------------------------------------

\draw[fwd, blue!70]  (hostA)   -- (pressA);
\draw[fwd, blue!70]  (pressA)  -- (signalA);
\draw[fwd, green!60!black] (signalA) -- (outA)
  node[midway, above, font=\scriptsize, text=green!50!black]
  {IVF aligns};

%% ---------------------------------------------------------------
%% Arrows — bottom path
%% ---------------------------------------------------------------

\draw[fwd, red!70]   (hostB)   -- (pressB);
\draw[fwd, red!70]   (pressB)  -- (signalB);
\draw[fwd, orange!80!black] (signalB) -- (outB)
  node[midway, above, font=\scriptsize, text=orange!60!black]
  {IVF drifts};

%% ---------------------------------------------------------------
%% Dashed arrows from IVF node to both paths
%% ---------------------------------------------------------------

\draw[dashed_arr] (ivf.north) -- (pressA.south)
  node[midway, left, font=\scriptsize, text=orange!70!black]
  {intensifies};

\draw[dashed_arr] (ivf.south) -- (pressB.north)
  node[midway, left, font=\scriptsize, text=orange!70!black]
  {intensifies};

%% ---------------------------------------------------------------
%% Labels at far left (hypothesis annotation)
%% ---------------------------------------------------------------

\node[font=\scriptsize\itshape, text=blue!70, left=0.5cm of hostA,
      align=right, anchor=east]
  {\textbf{H:} directional\\feedback};

\node[font=\scriptsize\itshape, text=red!70, left=0.5cm of hostB,
      align=right, anchor=east]
  {\textbf{H:} isotropic\\feedback};

\end{tikzpicture}
